% Part: model-theory
% Chapter: basics
% Section: isomorphism

\documentclass[../../../include/open-logic-section]{subfiles}

\begin{document}

\olfileid{mod}{bas}{iso}
\olsection{بنى متشاكلة}

تشابه !!{structure}s الرتبة الأولى على وجهين. فقد تتفق في
ال!!{sentence}s التي تصدق فيها؛ وتسمى الـ!!{structure}s عندئذ
\emph{متكافئة أوليًا}. وليس هذا الاتفاق مانعًا من اختلاف كبير
بينها: فقد تصدق فيها الـ!!{sentence}s نفسها، وتكون إحداها
!!{enumerable} والأخرى غير معدودة. فإن للـ!!a{structure} خصائص
لا تعبّر عنها لغة الرتبة الأولى، إما لقلة ما في اللغة من رموز،
وإما لحدود تلازم منطق الرتبة الأولى، كما تدل مبرهنة لوفنهايم--سكولم.
وأما الوجه الآخر لتشابه الـ!!{structure}s فأشد: أن تتحد صورتها
البنيوية، ولا تختلف إلا في الأشياء التي تتألف منها. وهذا هو
المعنى الذي نضبطه بمفهوم \emph{التشاكل}.

\begin{defn}
\ollabel{defn:elem-equiv}
لتكن !!{structure}s اثنتان $\Struct{M}$ و$\Struct M'$ للـ!!{language}
نفسها~$\Lang{L}$. نقول إن $\Struct{M}$ \emph{مكافئة أوليًا لـ}
$\Struct M'$، ونكتب $\Struct{M} \equiv \Struct M'$، إذا وفقط إذا كان،
لكل !!{sentence}~$!A$ في~$\Lang{L}$، $\Sat{{\OLId{latin-looped}{M}}}{!A}$ إذا وفقط إذا كان
$\Sat{{\OLId{latin-looped}{M}}'}{!A}$.
\end{defn}

\begin{defn}
\ollabel{defn:isomorphism} لتكن !!{structure}s اثنتان $\Struct{M}$
و$\Struct M'$ للـ!!{language} نفسها~$\Lang L$. نقول إن $\Struct{M}$
\emph{متشاكلة مع}~$\Struct M'$، ونكتب $\Struct{M} \simeq \Struct M'$،
إذا وفقط إذا وجدت دالة ${\OLId{latin-ordinary}{h}} \colon \Domain{{\OLId{latin-looped}{M}}} \to \Domain{{\OLId{latin-looped}{M}}'}$ بحيث:
\begin{enumerate}
\item تكون ${\OLId{latin-ordinary}{h}}$~!!{injective}: إذا كان ${\OLId{latin-ordinary}{h}}({\OLId{latin-ordinary}{x}}) = {\OLId{latin-ordinary}{h}}({\OLId{latin-ordinary}{y}})$، فإن ${\OLId{latin-ordinary}{x}} = {\OLId{latin-ordinary}{y}}$؛
\item تكون ${\OLId{latin-ordinary}{h}}$~!!{surjective}: لكل ${\OLId{latin-ordinary}{y}} \in \Domain{{\OLId{latin-looped}{M}}'}$ يوجد
  ${\OLId{latin-ordinary}{x}} \in \Domain{{\OLId{latin-looped}{M}}}$ بحيث ${\OLId{latin-ordinary}{h}}({\OLId{latin-ordinary}{x}}) = {\OLId{latin-ordinary}{y}}$؛
\item \ollabel{defn:iso-const}لكل !!{constant}~${\OLId{latin-ordinary}{c}}$:
  ${\OLId{latin-ordinary}{h}}(\Assign{{\OLId{latin-ordinary}{c}}}{{\OLId{latin-looped}{M}}}) = \Assign{{\OLId{latin-ordinary}{c}}}{{\OLId{latin-looped}{M}}'}$؛
\item \ollabel{defn:iso-pred}لكل !!{predicate}~${\OLId{latin-looped}{P}}$ ذي ${\OLId{latin-ordinary}{n}}$ مواضع:
  \[
  \tuple{{\OLId{latin-ordinary}{a}}_{\OLMathNumeral{1}}, \dots, {\OLId{latin-ordinary}{a}}_{\OLId{latin-ordinary}{n}}}\in \Assign{{\OLId{latin-looped}{P}}}{{\OLId{latin-looped}{M}}} \quad\text{إذا وفقط إذا}\quad
  \tuple{{\OLId{latin-ordinary}{h}}({\OLId{latin-ordinary}{a}}_{\OLMathNumeral{1}}), \dots, {\OLId{latin-ordinary}{h}}({\OLId{latin-ordinary}{a}}_{\OLId{latin-ordinary}{n}})} \in \Assign{{\OLId{latin-looped}{P}}}{{\OLId{latin-looped}{M}}'};
  \]
\item \ollabel{defn:iso-func}لكل !!{function}~${\OLId{latin-ordinary}{f}}$ ذات ${\OLId{latin-ordinary}{n}}$ مواضع:
  \[
  {\OLId{latin-ordinary}{h}}(\Assign{{\OLId{latin-ordinary}{f}}}{{\OLId{latin-looped}{M}}}({\OLId{latin-ordinary}{a}}_{\OLMathNumeral{1}}, \dots, {\OLId{latin-ordinary}{a}}_{\OLId{latin-ordinary}{n}})) =
  \Assign{{\OLId{latin-ordinary}{f}}}{{\OLId{latin-looped}{M}}'}({\OLId{latin-ordinary}{h}}({\OLId{latin-ordinary}{a}}_{\OLMathNumeral{1}}), \dots, {\OLId{latin-ordinary}{h}}({\OLId{latin-ordinary}{a}}_{\OLId{latin-ordinary}{n}})).
  \]
\end{enumerate}
\end{defn}

\begin{thm}
\ollabel{thm:isom}
إذا كان $\Struct{M} \iso \Struct M'$، فإن $\Struct{M} \elemequiv
\Struct{M'}$.
\end{thm}

\begin{proof}
لنأخذ ${\OLId{latin-ordinary}{h}}$ تشاكلًا من $\Struct{M}$ على $\Struct M'$. إذا كان
لدينا تعيين~${\OLId{latin-ordinary}{s}}$، فـ ${\OLId{latin-ordinary}{h}} \circ {\OLId{latin-ordinary}{s}}$ تركيب ${\OLId{latin-ordinary}{h}}$ و${\OLId{latin-ordinary}{s}}$، أي التعيين
في $\Struct{M'}$ الذي يحقق $({\OLId{latin-ordinary}{h}} \circ {\OLId{latin-ordinary}{s}})({\OLId{latin-ordinary}{x}}) = {\OLId{latin-ordinary}{h}}({\OLId{latin-ordinary}{s}}({\OLId{latin-ordinary}{x}}))$.
ونثبت بالاستقراء على ${\OLId{latin-ordinary}{t}}$ و$!A$ دعويين أقوى من المطلوب:
\begin{enumerate}
  \item[a.] ${\OLId{latin-ordinary}{h}}(\Value{{\OLId{latin-ordinary}{t}}}{{\OLId{latin-looped}{M}}}[{\OLId{latin-ordinary}{s}}]) = \Value{{\OLId{latin-ordinary}{t}}}{{\OLId{latin-looped}{M}}'}[{\OLId{latin-ordinary}{h}}\circ {\OLId{latin-ordinary}{s}}]$.
  \item[b.] $\Sat{{\OLId{latin-looped}{M}}}{!A}[{\OLId{latin-ordinary}{s}}]$ إذا وفقط إذا كان
    $\Sat{{\OLId{latin-looped}{M}}'}{!A}[{\OLId{latin-ordinary}{h}} \circ {\OLId{latin-ordinary}{s}}]$.
\end{enumerate}
أما الأولى، فبرهانها استقراء على تعقيد~${\OLId{latin-ordinary}{t}}$، وتفصيله كما يأتي.
\begin{enumerate}
\item إذا كان ${\OLId{latin-ordinary}{t}} \ident {\OLId{latin-ordinary}{c}}$، فإن $\Value{{\OLId{latin-ordinary}{c}}}{{\OLId{latin-looped}{M}}}[{\OLId{latin-ordinary}{s}}] = \Assign{{\OLId{latin-ordinary}{c}}}{{\OLId{latin-looped}{M}}}$
  و$\Value{{\OLId{latin-ordinary}{c}}}{{\OLId{latin-looped}{M}}'}[{\OLId{latin-ordinary}{h}} \circ {\OLId{latin-ordinary}{s}}] = \Assign{{\OLId{latin-ordinary}{c}}}{{\OLId{latin-looped}{M}}'}$. ولذلك
  ${\OLId{latin-ordinary}{h}}(\Value{{\OLId{latin-ordinary}{t}}}{{\OLId{latin-looped}{M}}}[{\OLId{latin-ordinary}{s}}]) = {\OLId{latin-ordinary}{h}}(\Assign{{\OLId{latin-ordinary}{c}}}{{\OLId{latin-looped}{M}}}) = \Assign{{\OLId{latin-ordinary}{c}}}{{\OLId{latin-looped}{M}}'}$ (بالبند
  \olref{defn:iso-const} من \olref{defn:isomorphism}) $=
  \Value{{\OLId{latin-ordinary}{t}}}{{\OLId{latin-looped}{M}}'}[{\OLId{latin-ordinary}{h}} \circ {\OLId{latin-ordinary}{s}}]$.
\item إذا كان ${\OLId{latin-ordinary}{t}} \ident {\OLId{latin-ordinary}{x}}$، فإن $\Value{{\OLId{latin-ordinary}{x}}}{{\OLId{latin-looped}{M}}}[{\OLId{latin-ordinary}{s}}] = {\OLId{latin-ordinary}{s}}({\OLId{latin-ordinary}{x}})$ و
  $\Value{{\OLId{latin-ordinary}{x}}}{{\OLId{latin-looped}{M}}'}[{\OLId{latin-ordinary}{h}} \circ {\OLId{latin-ordinary}{s}}] = {\OLId{latin-ordinary}{h}}({\OLId{latin-ordinary}{s}}({\OLId{latin-ordinary}{x}}))$. ومن ثم
  ${\OLId{latin-ordinary}{h}}(\Value{{\OLId{latin-ordinary}{x}}}{{\OLId{latin-looped}{M}}}[{\OLId{latin-ordinary}{s}}]) = {\OLId{latin-ordinary}{h}}({\OLId{latin-ordinary}{s}}({\OLId{latin-ordinary}{x}})) = \Value{{\OLId{latin-ordinary}{x}}}{{\OLId{latin-looped}{M}}'}[{\OLId{latin-ordinary}{h}} \circ {\OLId{latin-ordinary}{s}}]$.
\item إذا كان ${\OLId{latin-ordinary}{t}} \ident {\OLId{latin-ordinary}{f}}({\OLId{latin-ordinary}{t}}_{\OLMathNumeral{1}}, \dots, {\OLId{latin-ordinary}{t}}_{\OLId{latin-ordinary}{n}})$، فإن
  \begin{align*}
    \Value{{\OLId{latin-ordinary}{t}}}{{\OLId{latin-looped}{M}}}[{\OLId{latin-ordinary}{s}}] & = \Assign{{\OLId{latin-ordinary}{f}}}{{\OLId{latin-looped}{M}}}(\Value{{\OLId{latin-ordinary}{t}}_{\OLMathNumeral{1}}}{{\OLId{latin-looped}{M}}}[{\OLId{latin-ordinary}{s}}], \dots, \Value{{\OLId{latin-ordinary}{t}}_{\OLId{latin-ordinary}{n}}}{{\OLId{latin-looped}{M}}}[{\OLId{latin-ordinary}{s}}]) \quad\text{و}\\
  \Value{{\OLId{latin-ordinary}{t}}}{{\OLId{latin-looped}{M}}'}[{\OLId{latin-ordinary}{h}} \circ {\OLId{latin-ordinary}{s}}] & = \Assign{{\OLId{latin-ordinary}{f}}}{{\OLId{latin-looped}{M}}'}(\Value{{\OLId{latin-ordinary}{t}}_{\OLMathNumeral{1}}}{{\OLId{latin-looped}{M}}'}[{\OLId{latin-ordinary}{h}} \circ
    {\OLId{latin-ordinary}{s}}], \dots, \Value{{\OLId{latin-ordinary}{t}}_{\OLId{latin-ordinary}{n}}}{{\OLId{latin-looped}{M}}'}[{\OLId{latin-ordinary}{h}} \circ {\OLId{latin-ordinary}{s}}]).
  \end{align*}
  وبفرض الاستقراء نعلم، لكل ${\OLId{latin-ordinary}{i}}$، أن
  ${\OLId{latin-ordinary}{h}}(\Value{{\OLId{latin-ordinary}{t}}_{\OLId{latin-ordinary}{i}}}{{\OLId{latin-looped}{M}}}[{\OLId{latin-ordinary}{s}}]) = \Value{{\OLId{latin-ordinary}{t}}_{\OLId{latin-ordinary}{i}}}{{\OLId{latin-looped}{M}}'}[{\OLId{latin-ordinary}{h}}\circ {\OLId{latin-ordinary}{s}}]$. لذلك
  \begin{align}
    {\OLId{latin-ordinary}{h}}(\Value{{\OLId{latin-ordinary}{t}}}{{\OLId{latin-looped}{M}}}[{\OLId{latin-ordinary}{s}}])
    & = {\OLId{latin-ordinary}{h}}(\Assign{{\OLId{latin-ordinary}{f}}}{{\OLId{latin-looped}{M}}}(\Value{{\OLId{latin-ordinary}{t}}_{\OLMathNumeral{1}}}{{\OLId{latin-looped}{M}}}[{\OLId{latin-ordinary}{s}}], \dots, \Value{{\OLId{latin-ordinary}{t}}_{\OLId{latin-ordinary}{n}}}{{\OLId{latin-looped}{M}}}[{\OLId{latin-ordinary}{s}}])) \notag\\
    & = \Assign{{\OLId{latin-ordinary}{f}}}{{\OLId{latin-looped}{M}}'}({\OLId{latin-ordinary}{h}}(\Value{{\OLId{latin-ordinary}{t}}_{\OLMathNumeral{1}}}{{\OLId{latin-looped}{M}}}[{\OLId{latin-ordinary}{s}}]), \dots,
    {\OLId{latin-ordinary}{h}}(\Value{{\OLId{latin-ordinary}{t}}_{\OLId{latin-ordinary}{n}}}{{\OLId{latin-looped}{M}}}[{\OLId{latin-ordinary}{s}}])) \ollabel{iso-1}\\
    & = \Assign{{\OLId{latin-ordinary}{f}}}{{\OLId{latin-looped}{M}}'}(\Value{{\OLId{latin-ordinary}{t}}_{\OLMathNumeral{1}}}{{\OLId{latin-looped}{M}}'}[{\OLId{latin-ordinary}{h}} \circ {\OLId{latin-ordinary}{s}}], \dots,
    \Value{{\OLId{latin-ordinary}{t}}_{\OLId{latin-ordinary}{n}}}{{\OLId{latin-looped}{M}}'}[{\OLId{latin-ordinary}{h}} \circ {\OLId{latin-ordinary}{s}}]) \ollabel{iso-2}\\
    & = \Value{{\OLId{latin-ordinary}{t}}}{{\OLId{latin-looped}{M}}'}[{\OLId{latin-ordinary}{h}}\circ {\OLId{latin-ordinary}{s}}] \notag
  \end{align}
  أما \olref{iso-1} فمقتضى البند \olref{defn:iso-func} من
  \olref{defn:isomorphism}، وأما \olref{iso-2} فمقتضى فرض الاستقراء.
\end{enumerate}
وأما الجزء (b) فنترك برهانه للقارئ تمرينًا.

وأما إذا كانت $!A$ جملة، فلا يتوقف صدقها على التعيينين~${\OLId{latin-ordinary}{s}}$
و${\OLId{latin-ordinary}{h}} \circ {\OLId{latin-ordinary}{s}}$، فيلزم
$\Sat{{\OLId{latin-looped}{M}}}{!A}$ إذا وفقط إذا كان $\Sat{{\OLId{latin-looped}{M}}'}{!A}$.
\end{proof}

\begin{prob}
أتمّ برهان الجزء (b) من \olref[mod][bas][iso]{thm:isom} مفصلًا،
وعيّن أين استُعملت كل خاصية من الخواص الخمس المحددة للتشاكل
في \olref[mod][bas][iso]{defn:isomorphism}.
\end{prob}

\begin{defn}
\emph{التشاكل الذاتي} لبنية~$\Struct{M}$ هو تشاكل من
$\Struct{M}$ على نفسها.
\end{defn}

\begin{prob}
لتكن بنية~$\Struct{M}$، ولتكن ${\OLId{latin-looped}{X}}$ مجموعة جزئية قابلة
للتعريف في $\Struct{M}$، وليكن ${\OLId{latin-ordinary}{h}}$ تشاكلًا ذاتيًا لـ$\Struct{M}$.
أثبت أن ${\OLId{latin-looped}{X}} = \Setabs{{\OLId{latin-ordinary}{h}}({\OLId{latin-ordinary}{x}})}{{\OLId{latin-ordinary}{x}} \in {\OLId{latin-looped}{X}}}$، أي إن ${\OLId{latin-looped}{X}}$ ثابتة تحت ${\OLId{latin-ordinary}{h}}$.
\end{prob}


\end{document}
